\chapter{Burn Treatment}

\section*{Definition and Overview}
Burn treatment is the set of first aid, medical, and surgical measures used to care for people who have sustained burns. Burns are injuries to the skin and deeper tissues caused by heat, chemicals, electricity, friction, or radiation. The management approach depends on the depth and size of the burn, its location, the cause, and the person's age and general health.

Minor burns can usually be managed at home or by a family doctor, whereas deep, large, or critical burns are serious injuries requiring urgent specialist care, often in a dedicated burns unit. Early and correct first aid -- stopping the burning process and cooling the affected area -- can substantially reduce the severity of the injury, limit tissue damage, and improve the eventual cosmetic and functional outcome.

\section*{Epidemiology}
Burns are among the most common injuries worldwide and are largely preventable. Scalds from hot drinks, cooking, and bathing are the most frequent cause, particularly in young children and older adults. Men are more often injured through occupational fires, flame, and electrical injury. In Australia, thousands of people are admitted to hospital with burns each year, and globally burns are a leading cause of injury-related death and disability, with the highest burden in low- and middle-income countries. The vast majority of burns, however, are minor and never reach hospital care.

\section*{Aetiology and Pathophysiology}
Burns are classified by their cause and by their depth:
\begin{itemize}
    \item \textbf{Thermal burns:} flames, hot liquids (scalds), steam, and contact with hot objects; heat denatures proteins in the skin and kills the affected tissue
    \item \textbf{Chemical burns:} strong acids or alkalis in household or industrial products that continue to damage tissue until they are washed off
    \item \textbf{Electrical burns:} electric current heats the tissues and may cause deep damage, including heart rhythm disturbances and muscle injury out of proportion to the visible skin wound
    \item \textbf{Friction and radiation burns:} caused by rubbing surfaces or by sunlight and other radiation sources
\end{itemize}

Depth is conventionally described in degrees:
\begin{itemize}
    \item \textbf{Superficial (first degree):} damage limited to the outer layer of skin; red, painful, and dry, healing within days without scarring
    \item \textbf{Partial-thickness (second degree):} damage through to the deeper dermis; blistered, moist, and very painful, usually healing within one to three weeks
    \item \textbf{Full-thickness (third degree):} damage through the full thickness of the skin and sometimes deeper tissues; white, waxy, or charred, often relatively painless because nerve endings are destroyed, and generally requiring grafting
\end{itemize}

The larger and deeper the burn, the greater the risk of fluid loss, infection, and the generalised illness known as burn shock.

\section*{Clinical Features}
\begin{itemize}
    \item \textbf{Redness, pain, and tenderness} over the burned area, typical of superficial burns
    \item \textbf{Blisters}, which may be intact or broken, in partial-thickness burns
    \item \textbf{White, waxy, leathery, or charred skin} that may be painless in full-thickness burns
    \item \textbf{Swelling} of the injured area, which increases over the first few hours
    \item \textbf{Fluid loss:} weeping from broken blisters and, in large burns, signs of dehydration
    \item \textbf{Systemic features in severe burns:} shivering, low blood pressure, confusion, and reduced urine output
    \item \textbf{Inhalation injury:} singed nasal hairs, coughing, soot in the mouth, a hoarse voice, or breathing difficulty after fire or smoke exposure
\end{itemize}

\section*{Differential Diagnosis}
\begin{itemize}
    \item Superficial versus partial-thickness versus full-thickness burn; depth can be difficult to judge initially and may evolve over 24 to 48 hours
    \item Chemical versus thermal injury, which changes the approach to decontamination and treatment
    \item Electrical injury with internal organ damage versus an isolated skin burn
    \item Contact dermatitis, allergic skin reactions, or friction blisters that mimic a burn
    \item Blistering skin diseases, such as bullous disorders, occasionally mistaken for burns
\end{itemize}

\section*{When to Seek Medical Care}
For any burn, stop the burning process, cool the area with cool running water for 20 minutes, remove jewellery and tight clothing unless they are stuck to the skin, and cover the burn with a clean, non-stick dressing or cling film. Do not apply ice, butter, creams, or ointments.

\textbf{Call 000 (or local emergency services) for:}
\begin{itemize}
    \item Burns larger than the palm of the person's hand, deep burns, or burns involving the face, hands, feet, genitals, or a major joint
    \item Full-thickness burns with white, waxy, or charred skin
    \item Burns caused by electricity, chemicals, or explosions
    \item Burns associated with breathing difficulty, soot in the mouth or nose, or a hoarse voice
    \item Burns in young children, older adults, or people with significant pre-existing illness
    \item Any burn accompanied by signs of shock, such as pale, clammy skin, confusion, or fainting
\end{itemize}

\section*{Diagnosis and Investigations}
\begin{itemize}
    \item \textbf{History:} the cause of the burn, when it occurred, the mechanism of injury, and whether smoke or fumes were inhaled
    \item \textbf{Size estimation:} the burned area is calculated as a percentage of total body surface area, often using the rule of nines in adults
    \item \textbf{Depth assessment:} examination of colour, blistering, sensation, and capillary refill
    \item \textbf{Airway and breathing assessment:} essential in any facial burn or after fire exposure, to detect inhalation injury and airway oedema
    \item \textbf{Blood tests:} for significant burns, including full blood count, electrolytes, and kidney function
    \item \textbf{Monitoring:} urine output, heart rate, blood pressure, and oxygen saturation in hospital for clinically significant burns
\end{itemize}

\section*{Management}
\begin{itemize}
    \item \textbf{First aid:} stop the burning, cool with running water for 20 minutes, and cover with a clean dressing; blisters should not be burst and ice, butter, or creams avoided
    \item \textbf{Minor burns:} keep clean, apply a simple non-stick dressing, and use simple analgesia; intact blisters are usually left undisturbed
    \item \textbf{Wound cleaning:} the burn is cleaned, loose dead skin removed, and a non-stick dressing applied, sometimes with silver-based creams
    \item \textbf{Analgesia:} paracetamol or non-steroidal anti-inflammatory medicines for mild pain and stronger opioids for larger burns
    \item \textbf{Intravenous fluids:} for large burns, to replace lost fluid and protect organ function
    \item \textbf{Infection control:} regular dressing changes and antiseptic care, with antibiotics only if infection develops
    \item \textbf{Surgery:} debridement of dead tissue and skin grafting for full-thickness or deep partial-thickness burns
    \item \textbf{Specialist referral:} transfer to a burns unit for large, deep, or critical burns, and specialist follow-up for scar management
    \item \textbf{Rehabilitation:} physiotherapy and occupational therapy to maintain joint mobility and function, with pressure garments and silicone gels for scarring
\end{itemize}

\section*{Complications}
\begin{itemize}
    \item Wound infection, including bloodstream infection and sepsis
    \item Fluid and electrolyte disturbances from fluid loss through damaged skin
    \item Scarring, including raised hypertrophic scars and keloids
    \item Contractures, in which tightening of the skin and joints limits movement and may require surgery
    \item Chronic itch, altered sensation, or reduced sweating in healed areas
    \item Inhalation injury causing lung damage, pneumonia, and respiratory failure
    \item Psychological effects, including anxiety, depression, and post-traumatic stress
    \item In severe burns, multi-organ failure and death
\end{itemize}

\section*{Prognosis and Outcomes}
Most minor burns heal within one to three weeks with little or no permanent scarring. Deeper burns take longer, and the eventual outcome depends on good wound care, infection prevention, and appropriate surgical intervention. With modern specialist care, even many very large burns are survivable, but recovery can be prolonged and rehabilitation is essential for the best functional and cosmetic result. Scarring and stiffness may continue to improve for a year or more after the injury, and psychological support is often an important part of recovery.

\section*{Prevention and Self-Care}
\begin{itemize}
    \item Test bath water before bathing babies and children, and keep hot drinks and pot handles away from table edges
    \item Fit smoke alarms in the home, check them regularly, and practise a fire escape plan
    \item Set the hot water system to a safe temperature to prevent scalds
    \item Store matches, lighters, and chemicals out of children's reach
    \item Use protective clothing and safe practices when cooking, working with electricity, or handling chemicals
    \item Apply sunscreen and protect skin from excessive sun exposure
    \item After a burn, keep the wound clean, follow dressing instructions, and report any signs of infection promptly
\end{itemize}

\section*{Sources and Further Reading}
The following sources provide reliable, accessible information on burn treatment and first aid:
\begin{itemize}
    \item NHS. \textit{Burns and scalds.} https://www.nhs.uk/conditions/burns-and-scalds/
    \item World Health Organization. \textit{Burns fact sheet.} https://www.who.int/news-room/fact-sheets/
    \item Mayo Clinic. \textit{Burns first aid.} https://www.mayoclinic.org/first-aid/
    \item MedlinePlus. \textit{Burns.} https://medlineplus.gov/
    \item Centers for Disease Control and Prevention. \textit{Injury prevention and burns.} https://www.cdc.gov/
    \item MSD Manual. \textit{Burns.} https://www.msdmanuals.com/
\end{itemize}
